\documentclass[11pt]{article}
\usepackage[a4paper,margin=23mm]{geometry}
\usepackage[T1]{fontenc}
\usepackage{lmodern,amsmath,amssymb,amsthm,booktabs,longtable,array,microtype}
\usepackage{xcolor}
\usepackage[colorlinks=true,linkcolor=blue!50!black,urlcolor=blue!50!black]{hyperref}
\usepackage{enumitem}
\setlist{nosep,leftmargin=*}
\newtheorem{theorem}{Theorem}
\newtheorem{lemma}[theorem]{Lemma}
\newcommand{\R}[1]{\texttt{R12\_1\_#1}}
\newcommand{\FPdim}{\operatorname{FPdim}}
\renewcommand{\Vec}{\operatorname{Vec}}
\newcommand{\Rep}{\operatorname{Rep}}
\newcommand{\C}{\mathbb C}
\newcommand{\Z}{\mathbb Z}
\newcommand{\Q}{\mathbb Q}
\setlength{\parskip}{4pt}
\setlength{\parindent}{0pt}
\title{Rank twelve, multiplicity one:\\fresh audit and fourth theoretical continuation}
\author{Fusion-ring census research notes}
\date{Offline working report}
\begin{document}
\maketitle
\begin{abstract}
We audit all 438 exclusions and 148 unitary realizations retained in the
preceding round, and revisit its 46 unresolved based fusion rings. No
accepted classification is withdrawn. A complete Morita and extension-theory
argument resolves the remaining dimension-32 family, with two new unitary
models and sixteen exclusions. A new inner-conjugation equivariantization
argument excludes \R{332} over the complex numbers. A uniform integral-Lagrangian
criterion, together with complete quotient reconstructions, excludes all twenty
remaining dimension-48 rings. An alternate Brauer--Picard extension resolves
the four dimension-144 cases. The accepted totals are
478 excluded, 151 unitarily realized, and 3 unresolved, among 632 inputs.
Additional theoretical reductions for the unresolved families are stated
with their precise limitations. An erroneous blanket near-group assertion
in an earlier exploratory note is corrected; no accepted decision depended
on it. All work remains offline. The proofs and exact computations have
received independent checks within the assistant team; independent human
verification remains pending.
\end{abstract}

\section{Results, identifiers, and scope}
We count supplied \emph{based fusion rings}, not tensor-equivalence classes
of categories. Every input is bound to its complete multiplication tensor
and SHA256 digest. A negative decision excludes arbitrary complex fusion
categories. Every positive decision supplies an actual unitary category
and a unit-preserving bijection checking all $12^3=1728$ coefficients.

\begin{center}
\begin{tabular}{lrrr}\toprule
Stage & Excluded & Unitary models & Unresolved\\\midrule
Previous accepted ledger & 438&148&46\\
New decisions in this report &40&3&-43\\
Current accepted ledger &\textbf{478}&\textbf{151}&\textbf{3}\\\bottomrule
\end{tabular}
\end{center}
The new models are \R{320}, \R{419}, and \R{530}. ``New'' here means new
explicit realizations in this project; no novelty up to tensor equivalence
or exoticness is claimed. The new exclusions are:
\begin{center}\begin{tabular}{lp{126mm}}\toprule
Family & IDs, with prefix \texttt{R12\_1\_}\\\midrule
Dimension 32 &128, 144, 147, 173, 237, 275, 285, 294, 301, 318, 412, 415, 422, 431, 524, 536\\
Dimension 48 &130, 146, 151, 175, 216, 239, 244, 277, 289, 296, 303, 319, 407, 410, 421, 424, 433, 523, 526, 538\\
Dimension 144 &180,181,321\\
Quadratic &332\\\bottomrule\end{tabular}\end{center}

The proofs below separate categorical existence/converse arguments from
finite arithmetic and reconstruction checks. An associative candidate
table, a based-ring automorphism, or a failed construction search is not
treated as a categorical realization or exclusion.

\section{Audit of the preceding 586 settled decisions}
\subsection{Input, exclusions, and general-complex scope}
The fresh audit binds all 632 tensors to their original physical input
lines, verifies their fusion-ring axioms, and checks all 4012 recorded
closed-subring embeddings. Each of the 438 exclusion routes was reviewed
and replayed with explicit checks active under optimized Python.

Particular attention was paid to four possible sources of overstatement.
First, the published version of Schopieray's Lemma~3.1 gives a
\emph{pseudo-unitary Galois conjugate} for a two-FP-value categorification;
it does not assert pseudo-unitarity of the original category. Second,
unitary obstruction theorems were used only after a valid bridge to
unitarity: the relevant prime-power-product dimensions are solvable by
ENO and hence completely unitary by Galindo--Hong--Rowell. Third,
Liptrap negative certificates retain arbitrary nonzero complex variables,
including variable $\tau$; their integer Laurent relations cancel every
unknown, rather than rejecting a bounded phase search. Fourth, the
inherited rank-eight signed obstruction enumerates all 8192 unrestricted
oriented sign tables. All sixteen associative survivors are actual sign
gauges, and every possible spherical dimension-character orbit is
obstructed.

The twenty-four inherited exclusions are freshly mapped to six complete
lower-rank source proofs. The twelve field/order positivity checks were
strengthened to use rational isolating intervals. The dimension-120
premise checker was also strengthened: it now explicitly checks the
\emph{complete left and right stabilizer} of every dimension-two simple,
not just the common involution. The stronger hypothesis holds for all
sixteen inputs. The pointed Morita converse, finite-group restriction,
relative cohomology exhaustion, and coprime-twist lemma were separately
reread. All four resulting models were freshly reconstructed, and all
64 model/input comparisons were rechecked.

For \R{350}, the portable independent verifier reconstructs thirty low
induction frames, four PSD survivors, 45 complete joint cyclotomic trace
types, and 32 exact Farkas systems. It imports neither an optimizer nor
the earlier rank-six implementation. The codegree of each two-dimensional
matrix block is correctly $6$, not its Casimir eigenvalue $12$. The
cyclotomic trace box is not an assumed bound on the categorical exponent.

\subsection{Models, categorical actions, and unitary source tables}
The positive audit checks 381 ordinary first-round witnesses and 21
independently reconstructed Liptrap witnesses on 127 distinct rings,
then the specialized witnesses for the additional twenty-one rings.
It checks actual source categories, tensor actions, algebra objects,
complete simple-module lists, and full tensor identifications.
The checks distinguish a central lift from a commuting basis element.

A delicate inherited premise was strengthened constructively. For both
actual pseudo-PSU associator tables used in the rank-eight source
constructions, the audit finds an exact positive diagonal gauge. Every
nonzero coefficient satisfies
\[
 \frac{(F^{-1})_{fe}}{\overline{F_{ef}}}
 =\frac{w_{abe}w_{ecd}}{w_{bcf}w_{afd}},\qquad w_{ijk}>0.
\]
The weights are rational powers of $2$ and $1+\sqrt2$. Thus all 136
associator blocks become unitary. The two sets of 2688 pentagon identities
and the corresponding half-braiding hexagons are checked directly.
This avoids identifying a saved pentagon solution with a published
unitary example merely because their fusion rings agree.

\subsection{A genuine correction, with no classification change}
An earlier exploratory dimension-48 note called every near-group quotient
$\mathrm{NG}(G,2)$ with $|G|=8$ impossible. This omits the nonabelian
possibilities $D_8$ and $Q_8$, which do admit such categories. No accepted
exclusion used that assertion. The new analysis restores these branches.
Exact orbit-ring checks allow the $D_8$ branch for \R{216} and both
$D_8,Q_8$ branches for \R{239}; known models provide additional controls.
The analogous particular branches for \R{175}, \R{319}, \R{407}, and
\R{523} fail the complete nonabelian orbit-ring comparison. These are
restrictions on possible central lifts, not whole-ring exclusions.

\section{Complete resolution of the remaining dimension-32 family}
\begin{theorem}
For the verified input profile
$1^8,2^2,(\sqrt8)^2$, with adjoint series of dimensions $32,16,4,1$,
the only categorifiable supplied based rings are
\[
 \R{28},\ \R{183},\ \R{240},\ \R{405},\
 \R{408},\ \R{419},\ \R{522},\ \R{530}.
\]
All eight have unitary models. Consequently the other sixteen remaining
inputs in this family are excluded over $\C$.
\end{theorem}
This is a classification for the stated, exactly checked profile; it is
not a claim about every fusion category of dimension~32.

\subsection{A forced Morita reduction}
Let $\mathcal C$ be a hypothetical categorification. Its dimension-sixteen
integral part has eight invertibles and two dimension-two simples $X,Y$,
with common abelian stabilizer $H$ of order four. Write $Y=gX$ and
$P=\operatorname{Ad}_g|_H$. Liptrap's identities give the nondegenerate
bicharacter relations
\[
 B_Y(a,b)=B_X(Pa,b),\qquad B_{X^*}(a,b)=B_X(b,a),\qquad P^2=1.
\]
There is a nontrivial involution $h\in H$ fixed by $P$ and isotropic for
both forms. For $H=C_4$, take its unique involution. For $H=C_2^2$, a
symmetric nondegenerate form has an isotropic vector, and a nontrivial
orthogonal involution has an isotropic fixed line. If duality exchanges
$X,Y$, the cosquare relation forces $P=1$ in dimension two.

The rank-one action on $X$ trivializes the pointed associator on $H$.
Thus $A=1+h$ is an actual separable algebra. The isotropy identities make
the integral $A$-bimodule dual pointed of dimension sixteen. Its whole
Morita dual $\mathcal E$ has one of the two types
\[
 1^{16},(\sqrt8)^2,\qquad 1^{16},(\sqrt2)^8.
\]
The inverse algebra is an actual $1+c$, with $c$ a noncentral involution
in $G(\mathcal E)$, and the original integral type forces
$|C_{G(\mathcal E)}(c)|=8$.

\subsection{The two-square-root-eight source branch}
The common stabilizer has order eight. If the universal grading is $C_4$,
the associated involutive cosquare is the identity for every abelian group
of order eight; the invertible group would be abelian, a contradiction.
For universal group $C_2^2$, the inverse involution cannot lie in the
stabilizer. Indeed, for an orthogonal involution $P$ and a moved order-two
element $c$, the two required commutators
$B(c,c)=B(Pc,c)=-1$ are incompatible. For $C_2^3$ this follows from the
rank-one symmetric form associated to $P-1$; the $C_8$ and $C_4\times C_2$
cases have no suitable moved involution. Complete finite form checks
independently verify these elementary lemmas.

Therefore the source is a genuine crossed product of a TY category on an
abelian group of order eight. Its inverse algebra yields a coherent
$C_2$ equivariantization. Orthogonal involutions fixing four group elements
and the two parity-square choices exhaust the possibilities. They yield
only the already known rings \R{28}, \R{183}, \R{405}, \R{522}.

\subsection{The generalized-Ising sources and the missed nonsplit class}
For the other branch, $\mathcal E_{\rm ad}=\Vec_{C_2}$ and its universal
grading group $U$ has order sixteen. Put $K=G(\mathcal E)/\langle h\rangle$.
The inverse involution has a noncentral image $\bar c\in K$: a central
image would force a dimension-two object times an odd object to have
multiplicity two. Thus $K=D_8$. Write
\[
 G(\mathcal E)=\langle R,c,h\mid R^4=c^2=h^2=1,\ h\text{ central},
 \ cRc=R^{-1}h^b\rangle,\qquad b=0,1.
\]
The centralizer condition forces $R^4=1$. Both classes $b=0,1$ must be
retained; the nonsplit class $b=1$ is the new source of the two models.
The condition that no odd degree centralizes $\bar c$ leaves
$U=D_{16}$ or $SD_{16}$. There are therefore four reference families.

Their categorical exhaustion follows from ENO extension theory, not from
a search for associators. The coefficient module is
\[
 M=\operatorname{Inv}Z(\Vec_{C_2})=C_2^2
   =\operatorname{Ind}_K^U C_2,
 \qquad H^2(U,M)\cong H^2(K,C_2).
\]
The Shapiro projection is restriction to $K$ followed by forgetting the
half-braiding. It is exactly the central-extension class of the invertible
group displayed above, so no additional $H^2$ datum is omitted. Actual
Liptrap reference categories exist for all four possibilities; all others
are $U$-three-cocycle twists.

The LHS bounds for $|H^3(U,\C^\times)|$ are $32$ for $D_{16}$ and $16$
for $SD_{16}$. Explicit rotation, reflection, and (in the dihedral case)
mixed cocycles attain these bounds. The inverse $C_2$ algebra exists
precisely when the restriction is trivial; its algebra cochain is then
unique up to isomorphism. All rotation powers, admissible mixed twists,
and both inverse lifts are included.

The exact catalogue comprises 96 generalized-Ising constructions and 46
TY constructions. An independent implementation recomputes all 142 full
fusion tensors by monomial Hom equations. A separate graph-isomorphism
check proves that none of the sixteen excluded inputs is isomorphic to
an allowed output. The new nonsplit rotation twists give \R{419} and
\R{530}. Their actual source equations, Fourier unitarity, cocycle
identities, subgroup algebra, and complete bimodule tensors are also
audited separately.

\section{A new symmetry obstruction: \R{332}}
\begin{lemma}[Inner conjugation]
Let $H$ be a subgroup of invertible isomorphism classes of a fusion
category, preserving a fusion subcategory $\mathcal D$ by conjugation.
Then the conjugation functors give a coherent categorical $H$ action on
$\mathcal D$, even when the pointed associator on $H$ is nontrivial.
\end{lemma}
\begin{proof}
Choose representatives $u_h$ and multiplication isomorphisms
$m_{h,k}:u_h\otimes u_k\to u_{hk}$. Conjugation is
$A_u(X)=u\otimes X\otimes u^*$, with tensorator supplied by evaluation.
An isomorphism $f:u\to v$ induces
$A_f=f\otimes\mathrm{id}\otimes(f^{-1})^*$.
The two composites of multiplication maps on three invertibles differ by
the scalar associator. Conjugation applies this scalar and its inverse
on the two outer factors, so the discrepancy cancels. The induced
composition constraints therefore satisfy the action pentagon.
No algebra $\bigoplus_h u_h$ is asserted by this argument.
\end{proof}

\begin{theorem}
The fusion ring \R{332} has no complex categorification.
\end{theorem}
\begin{proof}
Pass to a pseudo-unitary Galois conjugate using published Schopieray
Lemma~3.1. The exact source contains a Haagerup-$C_3$ subcategory with
\[
 \alpha^3=1,\quad \rho\alpha=\alpha^{-1}\rho,\quad
 \rho^2=1+\rho+\alpha\rho+\alpha^2\rho,\quad
 d=\frac{3+\sqrt{13}}2.
\]
An existing invertible involution $b$ conjugates $\alpha$ to
$\alpha^{-1}$ and fixes $\rho$. The lemma gives a genuine $C_2$ action.
Its equivariantization is pseudo-unitary and has six simples
$(1,g,U,X,Y,V)$ of dimensions $(1,1,2,d,d,2d)$, with forgetful images
\[
 1,1,\alpha+\alpha^2,\rho,\rho,\alpha\rho+\alpha^2\rho.
\]
Here $g$ is the canonical sign object, $Y=gX$, and $gU=U,gV=V$.
The tensor images force $U^2=1+g+U$, $UX=UY=V$, and
$UV=X+Y+V$. Initially
$X^2=a+W+V$ with $a\in\{1,g\}$ and $W\in\{X,Y\}$.
If $a=g$, duality exchanges $X,Y$, but $Y^2=X^2$ by centrality of $g$,
contradicting the unique $X/Y$ summand. Thus $a=1$, and relabeling permits
$W=X$. All remaining products are forced:
\begin{align*}
 X^2=Y^2&=1+X+V,& XY&=g+Y+V,\\
 XV=YV&=U+X+Y+2V,&
 V^2&=1+g+U+2X+2Y+4V.
\end{align*}
This covers arbitrary action tensorators and coherence, not just one
chosen equivariant structure.

The global dimension is $D=39+9\sqrt{13}$. The character
\[
 (1,g,U,X,Y,V)\longmapsto(1,-1,0,\phi,-\phi,0),
 \qquad\phi=\frac{1+\sqrt5}{2},
\]
has formal codegree $f=5+\sqrt5$. Central induction would supply an
object of dimension $D/f$, an algebraic integer. But
\[
 N_{\Q(\sqrt{13},\sqrt5)/\Q}(D/f)
 =\frac{468^2}{20^2}=\frac{13689}{25}\notin\Z.
\]
This is impossible. The norm argument is independent of the project's
older exclusion of the isomorphic rank-six ring.
\end{proof}

\section{A uniform exclusion of the dimension-48 family}
\begin{theorem}
None of the twenty previously unresolved inputs of type $1^8,2^2,4^2$
listed above admits a complex categorification.
\end{theorem}
The proof has two independent parts: a categorical converse forcing any
hypothetical category to be group-theoretical, and complete exhaustion of
that group-theoretical construction. Neither a failed search nor an assumed
complete census of possible quotient rings is used.

\subsection{Complete group-theoretical parameters}
Double-coset dimensions and normal generation restrict the possible
finite-group/subgroup pairs. LHS and Kunneth bounds, attained by explicit
cocycles, give complete relative associator coverage for $S_4\times C_2$
and $A_4\rtimes C_4$. For $GL_2(3)$ and the binary octahedral group, the
sphere-bundle Gysin sequence and $p_1(\mathrm{Ad})=-4c_2$ identify the
periodic generators and their restriction kernels. All subgroup Schur
classes are retained. There are 432 exact restriction certificates.

The 92 admissible order-four/order-eight subgroup constructions give
36 wrong ranks, sixteen tables with a coefficient larger than one, and
forty copies of \R{30}, \R{414}, \R{417}, \R{531}. The remaining twelve
order-sixteen subgroup constructions give six wrong ranks and six copies
of \R{30} or \R{417}. Independent code certifies the complete projective
simple lists and reconstructs all relevant fusion tensors. Thus the twenty
targets have no group-theoretical realization.

\subsection{An integral-Lagrangian criterion}
\begin{lemma}
Every integral $2$-group extension of $\Vec_G^\omega$ is group-theoretical
when $G=A_4\times C_2$ or $SL(2,3)$, for arbitrary $\omega$.
\end{lemma}
\begin{proof}
A Lagrangian in $Z(\Vec_G^\omega)$ is described by a normal abelian
subgroup $N$ and an admissible invariant twisted alternating bicharacter.
Its corresponding module category has uniform normalized simple dimension
$\sqrt{|N|}$. Call it integral when $|N|$ is a square.

An integral Morita component preserves these Lagrangians. Indeed canonical
module dimensions are multiplicative under the balanced tensor functor
(ENO, Propositions 2.2 and 4.13). If the three module stabilizer orders
are $n_0,n,n'$, a product of two simples gives
$\sqrt{n_0}\sqrt n=k\sqrt{n'}$ with $k$ a positive integer.
When $n_0$ is a square, $n,n'$ have the same rational square class.

For $A_4\times C_2$, the square-order normal abelian subgroups are exactly
$1,V_4$. The first supplies one Lagrangian. The admissible forms on $V_4$
are either absent or a torsor under its two invariant alternating forms.
There are therefore one or three integral Lagrangians. For $SL(2,3)$,
only the trivial normal abelian subgroup has square order, giving one.
A finite $2$-group acting on an odd set fixes an element; the invariant
Lagrangian criterion then proves the claim.
\end{proof}
Integrality is preserved by the Morita operations used here. One direct
proof is that $Z(\mathcal C)$ is integral when $\mathcal C$ is integral.
For a Morita equivalent $\mathcal D$, induction gives
$\FPdim I_{\mathcal D}(X)=\FPdim(\mathcal D)\FPdim(X)\in\Z$.
Thus $\FPdim(X)$ is rational and an algebraic integer, hence an integer.

\subsection{All possible central quotients}
Each target has trivial universal grading. Solvability forces a final
$C_2$ equivariantization: a final extension is impossible, and a final
$C_3$ equivariantization would put an order-three subgroup among its eight
invertibles. Every possible central involution is retained. Its fixed
simple objects give four kinds of dimension-24 quotient.

\emph{Rank-twelve quotients.} A new bounded reconstruction starts with all
five groups of order eight, all eligible involutions, all left and right
actions on four dimension-two objects, all dualities, and all remaining
fusion coefficients. The invariant orbit products reduce the last step
to at most three binary variables. It finds exactly the two quotient
rings \R{351}, \R{513}; 704 full tables and every returned identification
were independently checked. This establishes completeness without relying
on the supplied rank-twelve list being exhaustive.

Both quotients are $C_3$ extensions of their pointed adjoint category of
dimension eight. For \R{351}, the group is $C_2^3$ and the distinct left
and right order-four stabilizers force its associator to be trivial:
the full bimodule stabilizer trivializes $\omega\boxtimes\omega^{-1}$,
and restriction to its order-two elements kills all seven Boolean
cohomology coordinates. An order-three isometry of the resulting
six-dimensional hyperbolic quadratic space has moving dimension two or
four. Moving dimension four has an invariant Lagrangian. Moving dimension
two forces the left and right stabilizers to coincide, contrary to the
actual tensor. Thus the category is group-theoretical. For \R{513},
the two cyclic stabilizers in $C_4\times C_2$ leave only the trivial
associator or the inflated $C_2$ class. Their centers have respectively
ten and four Lagrangians. A $C_3$ action fixes one.

The graded pointed Morita presentation has group $\Gamma=A\rtimes C_3$,
$|A|=8$. Nontrivial actions give only $A_4\times C_2$ or $SL(2,3)$.
A trivial action would split the category as its original pointed neutral
component times $\Vec_{C_3}$, contradicting the dimension-two objects.
Apply the integral-Lagrangian lemma to the Morita transport of the final
$C_2$ crossed product. Thus these quotient branches yield only
 group-theoretical parents.

\emph{Near-group quotients.} The restored $\mathrm{NG}(D_8,2)$ and
$\mathrm{NG}(Q_8,2)$ branches are treated in full. Solvability and complete
unitarizability justify the unitary classification. Each is Morita
equivalent to $\Vec_{S_4}^{\omega}$ with $\omega$ inflated from $S_3$.
Its center has exactly three Lagrangians, so every $C_2$ action fixes one.
Consequently every such equivariantization is group-theoretical.

\emph{Free rank-six quotients.} Their type is $1^4,2,4$ and their universal
grading is trivial. Solvability gives a further $C_2$ equivariantization
of dimension twelve. A complete 556-table elementary ring check, together
with the required swap of the two dimension-two objects and the
Jordan--Larson theorem, forces the predecessor ring $R_{3,V_4}$.
The outer action fixes the unique central lift of the inner boson; this
flattens the two equivariantizations into one by a group $P$ of order four.
There are three group-theoretical predecessors, with exactly three central
Lagrangians each, so their $P$ equivariantizations are group-theoretical.
For the three nongroup-theoretical predecessors, the exact soft-kernel
calculation gives an embedding of tensor autoequivalences into $S_3$.
An order-four action reduces to the trivial action or the reflection,
and in both cases a fusion coefficient two survives. Thus none produces
a target of multiplicity one. All categorical action coherences are
retained; the reduction is not a based-ring automorphism assumption.

\emph{Rank-nine quotients.} A separate direct reconstruction retains all
five dimension-two objects, their actions and dualities, and all bounded
integer coefficient assignments. Two of the four orbit-table classes
have no associative lift. Each of the other two has one lift, with
trivial universal grading and a unique central invertible involution $h$;
$h$ fixes all five noninvertibles. Solvability therefore forces an
$h$-equivariantization of a pointed dimension-twelve category. The outer
action fixes $h$, and its half-braiding is unique because the universal
grading is trivial. The central Tannakian subcategories combine to
$\Rep(P)$, $|P|=4$. The original category is a $P$ equivariantization
of that pointed category, and is group-theoretical. In particular this
argument does \emph{not} assume a $C_3$ grading for a rank-nine quotient;
some of the surviving based-ring lifts have no such grading.

These four branches cover all possible central involutions and contradict
the group-theoretical exhaustion above.

\section{Complete resolution of the dimension-144 family}
\begin{theorem}
\R{320} has a unitary model. The rings \R{180}, \R{181}, \R{321}
have no complex categorification.
\end{theorem}
The actual rank-three subcategory is $\Rep(S_3)$, so its regular algebra
$A=1+z+2E$ exists. The strengthened Morita argument proves that the
complete $A$-bimodule dual has type
\[
 1^{24},\ 2^{12},\ (\sqrt2)^{36},\qquad
 \mathcal D_{\rm ad}=\Vec_{V_4}.
\]
In particular the dimension-24 subcategory with $\Rep(S_4)$ fusion rules
has pointed dual $\Vec_{S_4}^{\omega}$, with $\omega|_{A_4}=0$.
It is never identified with ordinary $\Rep(S_4)$ just from its fusion ring.
The proof uses internal ends of transported modules, full stabilizers,
and the commuting left/right $S_4$ actions. These are spelled out in
\texttt{d144/new\_morita\_reduction.md}.

The universal group has order 36 and embeds in
$\mathrm{BrPic}(\Vec_{V_4})=(S_3\times S_3)\rtimes C_2$.
Here the quadratic space is the sum of two anisotropic planes, with
diagonal electric Lagrangian; its pointed parabolic is
$\operatorname{diag}(S_3)\times C_2$.
Only two index-two subgroups have the necessary six pointed grades:
\[
 U_{\rm base}=S_3\times S_3,\qquad
 U_{\rm alt}=\{(a,b,\epsilon):\operatorname{sgn}(a)
 \operatorname{sgn}(b)=(-1)^\epsilon\}\cong C_3^2\rtimes C_4.
\]
The twelve admissible marked base extensions are exhausted by the six
triality gaugings and their central zests. An explicit scalar cup cocycle
proves that these zests stay within the base subgroup; all give previously
modeled rings, none of the four targets.

For the alternate subgroup, put $M=\operatorname{Inv}Z(\Vec_{V_4})$.
The normal $C_3^2$ has no invariants in $M$, so $H^i(U_{\rm alt},M)=0$.
The obstruction $O_3$ vanishes and the multiplication datum is unique.
The obstruction $O_4$ restricts trivially to $C_3^2$, using the existing
base precursor, and to the Sylow $C_4$, whose fourth cohomology with
$\C^\times$ coefficients vanishes. Transfer kills $O_4$ entirely.
Actual alternate extensions therefore exist. Restriction of their
$H^3$ torsor onto the reverse $S_3$ subgroup is surjective, with kernel
of order two, so both admissible classes carry an actual separable
$S_3$ algebra. This proves existence before computing fusion rules.

The absolute three-primary parameter needs care. A Galois automorphism
fixing $\Q(\zeta_8)$ and inverting $\zeta_3$ acts affinely by $-1$ on
the $C_3^2$ extension torsor, so has a unique fixed point. The complete
marked first-axis triality table at $m=0$ is defined over $\Q(\zeta_8)$.
Its first-axis restriction and trivial reverse diagonal restriction make
its torsor point fixed. The alternate extension has the same property,
since its invariant line $x^2+y^2$ restricts injectively to the diagonal.
The remaining two-primary integral ambiguity is detected on the reverse
$S_3$ subgroup and vanishes because that restriction is fixed as trivial.
Hence its integral reverse fusion ring is exactly the $m=0$ ring.
This excludes \R{181}, \R{321}: their three dimension-four simples
have $N_{WW}^{W}=1$, whereas the forced ring has zero.

Finally the alternate Sylow $C_4$ restriction has a unique admissible
Brauer--Picard class up to pointed-parabolic conjugacy, and $H^2(C_4,M)=0$.
A concrete unitary Liptrap
source realizes it; its neutral $V_4$ associator is explicitly trivialized
by a two-cochain. Thus the local fusion tensor of the actual extension is
forced. The actual reverse $S_3$ bimodule orbits give four nonselfdual odd
objects of dimension $3\sqrt2$. This excludes \R{180}, which has two
self-dual odd objects. A short complete NIMrep argument then forces the
whole tensor: four of eight candidate tables are associative, and each
matches all 1728 coefficients of \R{320}. An independent implementation
reconstructs those tables. This closes the gap between precursor existence
and an actual model of the requested ring. Unitarizability follows from
solvability; no hypothetical lift of a ring automorphism is needed.

\section{The remaining quadratic cases and further reductions}
The unresolved quadratic rings are \R{561}, \R{623}, \R{628}.
The new inner-action method produces stronger finite section and
induction constraints. It must not be conflated with de-equivariantization
by a Tannakian subcategory. A braided central lift need not be Tannakian
in a nonbraided ambient fusion category. In particular the assertion
that every nondegenerate central pointed subcategory splits off as a
tensor factor is false in that generality. No exclusion in this report
uses that assertion. The partial constraints and the exact remaining
categorical gaps are retained in the evidence notes.

For \R{561} and \R{628}, the coherent inner $C_3$ action and complete
72-section arithmetic leave only subgroup sections. They imply an actual
central braided $C_3$ lift, but do not force its braiding to be isotropic.
The three surviving sections are compatible with the three possible
braided diagonals. In the nondegenerate cases, the actual regular
$C_3$ algebra gives an Azumaya Morita autoequivalence; it does not give a
rank-four tensor quotient or a different excluded ring.

For \R{623}, the corresponding inner $C_2$ reduction forces a central
semionic lift. Its square-root twists and duality greatly reduce the
possible central-induction data. Exact dimension and Gram constraints
leave thirteen integer profiles, and the strengthened semion trace
conditions eliminate one of them. Twelve remain unseparated. Neither a
Tannakian quotient nor an algebra $1+z$ is asserted for the semionic object.

As a separate uniform result, complete exact TY-center calculations give
Lagrangian counts $7,5,2,2,1,1,1,1$ for the eight dimension-eight
Tambara--Yamagami categories. Thus every $C_3$ extension of such a category
is group-theoretical. This useful result is not used to assign a false
$C_3$ grading to the surviving rank-nine quotient tables.

The September 2026 Haagerup--Izumi construction and classification were
checked against the remaining inputs. Their odd-order and stated
Q-system hypotheses do not cover these three even-order rings. The
paper's general inner-conjugation proposition independently supports the
action lemma used above, but supplies no missing classification here.

\section{Remaining IDs and reproducibility}
Exactly three rings remain unresolved: \R{561}, \R{623}, \R{628}.
No failed search is promoted to an exclusion. All other forty-three
formerly unresolved cases have either a general-complex obstruction or
an actual unitary realization.

The package contains the updated ledger, input hashes, full proof notes,
portable accepted constructions and verifiers, separate independent
receipts, and the frozen preceding package. Discovery failures and
conditional reductions are labeled separately from accepted decisions.
No repository commit, push, website edit, or deployment was made.

\section*{Selected primary references}
\begin{enumerate}
\item J.~Liptrap, \emph{Generalized Tambara--Yamagami categories},
\href{https://arxiv.org/abs/1002.3166}{arXiv:1002.3166}, Theorem4.
\item P.~Etingof, D.~Nikshych, V.~Ostrik, \emph{Fusion categories and
homotopy theory}, \href{https://arxiv.org/abs/0909.3140}{arXiv:0909.3140}.
\item A.~Schopieray, \emph{Categorification of integral group rings extended
by one dimension}, JLMS108 (2023), 1617--1641,
\href{https://doi.org/10.1112/jlms.12789}{published Lemma3.1}.
\item V.~Ostrik, \emph{Pivotal fusion categories of rank3},
\href{https://arxiv.org/abs/1309.4822}{arXiv:1309.4822}, Theorem2.13.
\item P.~Etingof, D.~Nikshych, V.~Ostrik, \emph{Weakly group-theoretical
and solvable fusion categories},
\href{https://arxiv.org/abs/0809.3031}{arXiv:0809.3031}, Theorem1.6.
\item C.~Galindo, S.-M.~Hong, E.~Rowell, \emph{Generalized and
quasi-localizations of braid group representations},
\href{https://arxiv.org/abs/1105.5048}{arXiv:1105.5048}, Theorem5.20.
\item M.~Izumi, H.~Tucker, \emph{Algebraic realization of noncommutative
near-group fusion categories},
\href{https://arxiv.org/abs/1908.01655}{arXiv:1908.01655}, Theorems1.2,3.2.
\item S.~Gelaki, D.~Naidu, D.~Nikshych, \emph{Centers of graded fusion
categories}, \href{https://arxiv.org/abs/0905.3117}{arXiv:0905.3117},
Corollaries 3.10--3.11.
\item D.~Nikshych, B.~Riepel, \emph{Categorical Lagrangian Grassmannians
and Brauer--Picard groups of pointed fusion categories},
\href{https://arxiv.org/abs/1309.5026}{arXiv:1309.5026}.
\item D.~Jordan, E.~Larson, \emph{On the classification of certain fusion
categories}, \href{https://ems.press/content/serial-article-files/30464}{Theorems 1.1 and 1.4}.
\item C.~Edie-Michell, \emph{Equivalences of graded categories},
\href{https://arxiv.org/abs/1711.00645}{arXiv:1711.00645}.
\item I.~Marshall, D.~Nikshych, \emph{On the Brauer--Picard groups of
fusion categories}, \href{https://arxiv.org/abs/1603.04318}{arXiv:1603.04318}.
\item T.~Gannon, A.~Schopieray, H.~Yadav, \emph{On Haagerup--Izumi fusion
categories}, \href{https://arxiv.org/abs/2609.25185}{arXiv:2609.25185},
Proposition6.1 and the explicit scope of TheoremA.
\end{enumerate}
\end{document}
